\subsection{The Dot Product}


\content{
\begin{definition} (The Dot Product) If $\vec{a}=\langle a_1,a_2,a_3\rangle$
and $\vec{b}=\langle b_1,b_2,b_3\rangle$, then the dot product
$\vec{a}\cdot\vec{b}$ is given by 
$$ \vec{a}\cdot \vec{b} = a_1b_1+a_2b_2+a_3b_3. $$
\end{definition}
\begin{example} if $\vec{a} =\langle 1, 0, 3\rangle$ and 
$\vec{b} =\langle 2,-1,1\rangle$, compute $\vec{a}\cdot \vec{b}$.
\end{example}
}{% presentation
\vspace{1in}
}{% nopresentation
}{% comments
}


\content{
\begin{theorem} (Properties of the Dot Product) if $\vec{a}$, $\vec{b}$, and 
$\vec{c}$ are vectors in $\R^3$ and $c$ is a scalar, then 
\begin{enumerate}
\item $\vec{a}\cdot \vec{a} = |\vec{a}|^2$
\item $\vec{a}\cdot\vec{b} = \vec{b}\cdot \vec{a}$
\item $\vec{a}\cdot(\vec{b}+\vec{c}) = \vec{a}\cdot \vec{b} +
\vec{a}\cdot\vec{c}$
\item $(c\vec{a})\cdot \vec{b} =c(\vec{a}\cdot \vec{b})=\vec{a}\cdot
(c\vec{b})$
\item $\vec{0}\cdot \vec{a} = 0$
\end{enumerate}
\end{theorem}
}{% presentation
\vspace{2in}
}{% nopresentation
}{% comments
Maybe provide a proof of 1 and 3.
}

\content{
\begin{theorem} If $\theta$ is the angle between vectors $\vec{a}$ and
$\vec{b}$, then 
$$ \vec{a}\cdot \vec{b} = |\vec{a}||\vec{b}|\cos\theta. $$
\end{theorem}
}{% presentation
\vspace{1in}
}{% nopresentation
}{% comments
In the space below, solve this equation for $\cos\theta$, since this is the
primary way that we will compute the angle between vectors. 
}

\content{
\begin{example} Find the angle between the vectors $\vec{a} = \la 2,2,-1 \ra$
and $\vec{b} = \la 5,-3,2\ra $.
\end{example}
}{% presentation
\vspace{2in}
}{% nopresentation
}{% comments
}

\content{
\begin{theorem} Two vectors are orthogonal (perpendicular) if their dot product
is zero. 
\end{theorem}
}{% presentation
\vspace{2in}
}{% nopresentation
}{% comments
Either prove this using the formula above, or use the Pythagorean Theorem. 
}

\subsubsection{Projections}

\content{
We will often need to break vectors up into components to determine how much
force or work is being done in a certain direction.  More specifically, if we
are given two vectors $\vec{a}$ and $\vec{b}$, we wish to determine the
projection of $\vec{a}$ onto $\vec{b}$. 
}{% presentation
\vspace{2in}
}{% nopresentation
}{% comments
Draw both vectors $\vec{a}$ and $\vec{b}$ from the same starting point, then
drop the perpendicular from $\vec{a}$ onto $\vec{b}$ and derive the projection
formula from there. 

\textbf{Derive both the component formula and the projection formula. }
}

\content{
\begin{example} Find the projection of the vector $\vec{a}=\la 1,1,2\ra $ onto
the vector $\vec{b}=\la -2,3,1\ra$.
\end{example}
}{% presentation
\vspace{2in}
}{% nopresentation
}{% comments
}

\content{
\begin{example} One use of projections is to find the amount of work done by a
constant force. A wagon is being pulled a distance of 100 m along by a constant
force of 70 N. The handle of the wagon is held at an angle of $35^\circ$ above
the horizontal. Find the work done by the force. 
\end{example}
}{% presentation
}{% nopresentation
}{% comments
Set up a diagram, then when the triangle is drawn, use the fact that work is
force times distance to derive the formula 
$$ W = \vec{F}\cdot \vec{D}$$ 
where $\vec{F}$ is the force vector being applied and $\vec{D}$ is the direction
with magnitude 100. 
}

\subsection{The Cross Product}

\content{
\begin{definition} (The Cross Product) If $\vec{a}=\la a_1,a_2,a_3\ra$ and 
$\vec{b} = \la b_1,b_2,b_3\ra$, then the cross product, denoted $\vec{a}\times
\vec{b}$ is given by 
$$ \vec{a}\times \vec{b} = \la a_2b_3-a_3b_2,a_3b_1 - a_1b_3,a_1b_2-a_2b_1\ra.$$
\end{definition}
}{% presentation
\vspace{2in}
}{% nopresentation
}{% comments
Show how this can be remembered by taking the determinant of a 3x3 matrix. 
}

\content{
\begin{example} Find the cross product between the vectors 
$\vec{a}=\la 1,3,4\ra$ and $\vec{b}=\la 2,7,-5\ra$.
\end{example}
}{% presentation
\vspace{2in}
}{% nopresentation
}{% comments
$\vec{a}\times\vec{b}=\la -43, 13, 1\ra$.
}

\content{
\begin{theorem} The vector $\vec{a}\times\vec{b}$ is orthogonal to both
$\vec{a}$ and $\vec{b}$.
\end{theorem}
}{% presentation
}{% nopresentation
}{% comments
}

\content{
\begin{theorem} If $\theta$ is the angle between vectors $\vec{a}$ and
$\vec{b}$, then 
$$ |\vec{a}\times\vec{b}| = |\vec{a}||\vec{b}|\sin\theta. $$
\end{theorem}
}{% presentation
\vspace{2in}
}{% nopresentation
}{% comments
Use the space below to show that this is precesily the area of the parallelogram
spanned by $\vec{a}$ and $\vec{b}$, so that the magnitude of the cross product
gives us the area of that parallelogram. 
}

\content{
\begin{theorem} Two vectors $\vec{a}$ and $\vec{b}$ are parallel if and only if 
$\vec{a}\times\vec{b}=0$. 
\end{theorem}
}{% presentation
}{% nopresentation
}{% comments
}

\content{
\begin{example} Find a vector which is perpendicular to the plane passing
through the points $P(1,4,6)$, $Q(-2,5,-1)$, and $R(1,-1,1)$.
\end{example}
}{% presentation
\vspace{3in}
}{% nopresentation
}{% comments
One answer is $\la -40, -15, 15\ra$.
}

\content{
\begin{theorem} (Properties of the Cross Product) If $\vec{a}$, $\vec{b}$, and 
$\vec{c}$ are vectors and $c$ is a scalar, then 
\begin{enumerate}
\item $\vec{a}\times\vec{b} = -\vec{b}\times\vec{a}$
\item $\vec{a}\times (\vec{b}+\vec{c}) =
\vec{a}\times\vec{b}+\vec{a}\times\vec{c}$
\item $(c\vec{a})\times\vec{b}=c(\vec{a}\times\vec{b}) = 
\vec{a}\times(c\vec{b})$
\item $(\vec{a}+\vec{b})\times \vec{c} = \vec{a}\times \vec{c} +
\vec{b}\times\vec{c}$
\item $\vec{a}\cdot (\vec{b}\times\vec{c}) = 
(\vec{a}\times\vec{b})\cdot\vec{c}$
\item $\vec{a}\times(\vec{b}\times\vec{c})=
(\vec{a}\cdot\vec{c})\vec{b}-(\vec{a}\cdot \vec{b})\vec{c}$
\end{enumerate}
\end{theorem}
}{% presentation
\vspace{4in}
}{% nopresentation
}{% comments
Maybe give a proof of one of these. 
}

\content{
The product $\vec{a}\cdot (\vec{b}\times\vec{c})$ that occurs in the above
theorem is known as the scalar triple product. And it's value is the volume of
the parallelpiped determined by the vectors $\vec{a}$, $\vec{b}$, and $\vec{c}$. 
}{% presentation
\vspace{2in}
}{% nopresentation
}{% comments
Don't give a proof, just draw out a parallelpiped, and move onto the next
example. 
}

\content{
\begin{example} Find the volue of the parallelpiped determined by the three
vectors $\vec{a}=\la 1,4,-7\ra$, $\vec{b}=\la 2,-1,4\ra$, and 
$\vec{c} =\la 0,-9,18\ra$. 
\end{example}
}{% presentation
\vspace{2in}
}{% nopresentation
}{% comments
The volume is 0, which means that these vectors actually lie in the same plane,
that is, they are coplanar. 

An Application to computing torque is outlined in the textbook for those who are
interested in reading that. 
}
